Scores that screen newly launched tokens are widely used and almost never validated. We audit the published DryDock score for Arc launches (v1), an additive 0 to 100 index of seven weighted factors with four hard caps {[}1{]}, and report twelve findings. The published weights sum to 90, not 100. The score has no stated estimand, so it cannot be tested. Market-activity factors carry 44\% of attainable points, although insiders can manufacture activity and it peaks when entry risk peaks. Insider concentration is thresholded on supply share alone, yet the price impact of an insider exit depends on that share relative to pool depth: under a constant-product pool a 15\% stake implies a 47\% or an 84\% price drop depending on the pool's share of supply. The ``profitable trader'' wallet label is awarded to 16\% to 32\% of no-skill wallets at five launches. Under a stated interpolation assumption, a launch whose measured insider share is 20\% still scores 76 (``Strong setup'') while that cluster's exit would cut a buyer's position by 89\%.

We then specify v2. It replaces points with three quantities that carry units. The \emph{structural floor} is the recovery of a \$100 reference position under the worst single structural event available at a given block, computed by simulation. \emph{Channel probabilities} are forecasts of three resolvable on-chain events (tradability loss, liquidity removal, insider exit), calibrated against outcomes. The \emph{Integrity Score} is 100 minus their expected loss. Insider holdings are reported as an identified interval whose pessimistic bound cannot be lowered by Sybil splitting. Price momentum is excluded from the score by construction. Forecasts are committed on-chain before outcomes and scored in public with proper scoring rules. We give five propositions, an eleven-attack analysis, a pre-registered validation protocol with sample sizes, and explicit ship and kill criteria.
